\documentclass[10pt,a4paper]{article}
\usepackage[margin=2.2cm]{geometry}
\usepackage{fontspec}
\usepackage{kotex}
\setmainfont{Apple SD Gothic Neo}
\setsansfont{Apple SD Gothic Neo}
\usepackage{amsmath,amssymb}
\usepackage{graphicx}
\usepackage{booktabs}
\usepackage{caption}
\usepackage[dvipsnames]{xcolor}
\usepackage[colorlinks=true,linkcolor=MidnightBlue,citecolor=MidnightBlue,
            urlcolor=MidnightBlue]{hyperref}
\usepackage{titlesec}
\titleformat{\section}{\large\bfseries}{\thesection}{0.6em}{}
\titleformat{\subsection}{\normalsize\bfseries}{\thesubsection}{0.5em}{}
\captionsetup{font=small,labelfont=bf}
\setlength{\parskip}{0.35em}
\setlength{\parindent}{0pt}

% 주장 상자 — 모든 주장은 반증 조건을 함께 적는다
\usepackage{mdframed}
\newmdenv[linewidth=0.4pt,linecolor=gray!60,backgroundcolor=gray!5,
          innertopmargin=6pt,innerbottommargin=6pt,skipabove=8pt,skipbelow=8pt]{claimbox}
\newcommand{\claim}[2]{\begin{claimbox}\textbf{주장.} #1\par\smallskip
  \textcolor{gray!70}{\textbf{반증 조건.} #2}\end{claimbox}}

\title{\bfseries 스무 노트 동안 89행이라 불렀는데 학습은 16행이었다}
\author{월드모델 실험실 · 노트 665}
\date{2026-08-06}
\begin{document}\maketitle

\section{한 줄}

이 세션이 스무 노트에 걸쳐 ``팝업 89행''을 시험대로 논했다. 그 89 중 모형이
\textbf{배우는 데 쓰는 것은 16행}이다. 유보가 59라 학습의 \textbf{3.7배}이고,
판 열한 도메인 중 \textbf{학습 행이 제일 적다}.

\section{어쩌다 봤나}

노트 664 가 잔여 물음을 남겼다 --- 시장팝업의 정밀 라벨(A·B, 54행)에서 방문자
축의 상관이 $-0.038$ 인데, 그것이 \emph{표본이 작아서}인지 \emph{정밀 라벨이
동네 유동인구와 덜 얽혀서}인지 안 쟀다.

팝업 도메인에도 등급이 있다(\verb|outcome.label_trust.grade| 78건, A 52 · B 16).
\textbf{등급 분포가 정반대라} 자연 실험이 된다 --- 시장은 기사 발표치(A 1 · B 53 ·
C 47), 팝업은 드라이브 결과보고서의 일별 표·계수기 로그(A 67\%).

\claim{팝업에서는 A·B 가 C·D·E 보다 높을 것이다 --- 팝업은 A 가 67\% 인데 도메인
전체가 $+0.019$ 이므로 그 이득이 A급에서 나와야 한다.}{팝업의 이득이 팝업의
학습 행에서 온 것이 아니면 이 추론은 성립하지 않는다.}

쟀더니 팝업의 축 관측 $\cdot$ 라벨 $\cdot$ 학습 구간을 다 만족하는 행이
\textbf{13}이었다. 사각으로 적어 둔 것(``C·D·E 가 10행뿐일 수 있다'')보다 훨씬
작아서 원인을 봤다.

\section{판의 학습/유보}

\begin{center}\small
\begin{tabular}{lrrr}\toprule
도메인 & 전체 & 학습($<$2025) & 유보(2025+) \\ \midrule
\textbf{팝업} & 75 & $\mathbf{16}$ & $\mathbf{59}$ \\
아이돌 & 81 & 54 & 25 \\
시장팝업 & 249 & 123 & 126 \\
도서 & 243 & 243 & 0 \\
게임 & 482 & 482 & 0 \\ \bottomrule
\end{tabular}\end{center}

팝업만 유보가 학습보다 크다. 다른 도메인은 반대이거나 유보가 아예 없다.

\section{그래서 노트 662 를 다시 읽어야 한다}

노트 662 가 방문자 축으로 팝업 도메인 $+0.019$ 를 얻었다. \textbf{그 이득은 팝업의
16 학습 행에서 나온 것일 수 없다.} 같은 열이 시장팝업의 123 학습 행에서 배운
것이 팝업 유보 59행에 옮겨 간 것이다 --- \textbf{도메인 안 학습이 아니라 L2
전이}다.

노트 660 이 ``열을 나눠 쓰면 한쪽이 벌고 한쪽이 낸다''로 적은 것의 기제가 이것이다.
그때는 덮음으로 설명했고 그것도 맞지만, \emph{누가 배우고 누가 쓰나}가 빠져 있었다.

\section{내 예측이 틀린 방식}

추론은 이랬다 --- ``팝업이 A 67\% 이고 도메인이 $+0.019$ 이니 A급에서 신호가
있어야 한다.'' \textbf{도메인 크기만 보고 학습/유보 비를 안 봤다.} 이득을 낸 것이
그 도메인의 학습 행이라고 암묵적으로 가정한 것이고, 팝업에서는 그것이 거짓이다.

\section{남은 물음은 그대로다}

합친 표본에서 무늬는 재현된다 --- \textbf{A·B $n{=}65$ $\rho$ $-0.0936$}
$[-0.332, +0.157]$ 대 \textbf{C·D·E $n{=}64$ $\rho$ $+0.3093$} $[+0.075, +0.525]$.
구간이 이번엔 거의 안 겹친다.

그런데 재현의 대부분이 시장팝업에서 온다(팝업 A·B 는 11행). \textbf{정밀 라벨
표본이 실질적으로 하나}이므로 원리적 한계인지 표본 문제인지 \emph{여전히 못
가른다}. KOPIS 가 필요한 이유가 셋이 됐다 --- 행 $\cdot$ A급 라벨 $\cdot$
\textbf{학습 행}.

\section{조항}

\claim{도메인 크기는 학습/유보를 갈라 적는다. 한 숫자로 적지 않는다.}{열값이
학습 행 수와 무관하거나 채점이 유보 행 수와 무관하면 이 조항이 불필요하다.}

열값은 \textbf{학습 행}이 정하고(모형이 그 열로 쪼개며 잡음을 배우는 곳)
채점은 \textbf{유보 행}이 정한다. 둘이 다른 물음에 답하므로 한 숫자로 적으면
어느 쪽 이야기인지 모른다. 노트 646 의 ``팝업은 판의 0.4\%''도 유보 기준이었다.

\begin{figure}[h]\centering
\includegraphics[width=\textwidth]{figs/f1.pdf}
\caption{(a) 팝업만 유보가 학습의 3.7배다. (b) 합친 표본에서 신호가 정밀 라벨이
아니라 하한값·추산 라벨에 있다. (c) 그래서 도메인 크기를 학습/유보로 갈라 적는다.}
\end{figure}

\end{document}
